\documentclass[11pt,a4paper]{article}
\usepackage[T1]{fontenc}
\usepackage[utf8]{inputenc}
\usepackage{lmodern}
\usepackage[margin=26mm,headheight=15pt]{geometry}
\usepackage{amsmath,amssymb,amsthm,mathtools,bm}
\usepackage{microtype,booktabs,array,longtable,enumitem,needspace}
\usepackage[dvipsnames]{xcolor}
\usepackage{xurl}
\usepackage[colorlinks=true,linkcolor=MidnightBlue,citecolor=MidnightBlue,urlcolor=MidnightBlue]{hyperref}
\usepackage{fancyhdr}
\pagestyle{fancy}\fancyhf{}
\fancyhead[L]{\small The lower critical endpoint}
\fancyhead[R]{\small Condensation and transseries}
\fancyfoot[C]{\thepage}
\renewcommand{\headrulewidth}{0.3pt}
\setlength{\parskip}{3pt}
\setlength{\parindent}{1.15em}
\setlist{itemsep=3pt,topsep=5pt}
\numberwithin{equation}{section}
\newtheorem{theorem}{Theorem}[section]
\newtheorem{lemma}[theorem]{Lemma}
\newtheorem{proposition}[theorem]{Proposition}
\newtheorem{corollary}[theorem]{Corollary}
\theoremstyle{definition}
\newtheorem{definition}[theorem]{Definition}
\newtheorem{example}[theorem]{Example}
\theoremstyle{remark}
\newtheorem{remark}[theorem]{Remark}
\newcommand{\R}{\mathbb R}
\newcommand{\C}{\mathbb C}
\newcommand{\N}{\mathbb N}
\newcommand{\E}{\mathbb E}
\newcommand{\Prob}{\mathbb P}
\newcommand{\ee}{\epsilon}
\newcommand{\dd}{\,\mathrm d}
\newcommand{\ind}{\mathbf1}
\newcommand{\dto}{\xrightarrow{\mathrm d}}
\newcommand{\pto}{\xrightarrow{\Prob}}
\newcommand{\law}{\mathcal L}
\newcommand{\tv}{d_{\mathrm{TV}}}
\newcommand{\coef}[2]{[#1]#2}
\DeclareMathOperator{\Li}{Li}
\DeclareMathOperator{\Var}{Var}
\DeclareMathOperator{\Pois}{Pois}
\DeclareMathOperator{\Rea}{Re}
\DeclareMathOperator{\Ima}{Im}
\DeclareMathOperator{\supp}{supp}
\newcommand{\commit}{\texttt{3d5973524506411392a911470b5ddc35521568ea}}
\hypersetup{pdftitle={The Lower Critical Endpoint: Arbitrary-Rate Condensation, Cauchy Cutoff Laws, and Degenerating Transseries},pdfauthor={Research article prepared for Vladimir Reshetnikov},pdfsubject={A model-specific answer to the exponent-one endpoint question in ProveIt}}
\begin{document}
\begin{center}
{\small\scshape Transseries, asymptotic inversion, and conditional probability}\par
\vspace{12pt}
{\LARGE\bfseries The Lower Critical Endpoint}\par
\vspace{8pt}
{\Large Arbitrary-Rate Condensation, Cauchy Cutoff Laws,\par and Degenerating Transseries}\par
\vspace{14pt}
{\normalsize Research article prepared for Vladimir Reshetnikov}\par
{\small 29 September 2026}\par
\end{center}
\vspace{6pt}
\begin{abstract}
We study the critical countable-action equation
\[
 U_\ee(q)=c_\ee\bigl(\Li_{2+\ee}(qe^{U_\ee(q)})+P(qe^{U_\ee(q)})\bigr),
 \qquad c_\ee=(\zeta(1+\ee)+P'(1))^{-1},
\]
where $P$ changes finitely many nonnegative action weights. For every joint
limit $n\to\infty$, $\ee=\ee_n\downarrow0$, with no restriction on its rate,
we prove
\[
 [q^n]U_\ee(q)\sim \rho_\ee^{-n}\frac{\ee}{n d_n},\qquad
 d_n=\left(\frac{nc_\ee}{\ee}\right)^{1/(1+\ee)},\quad
 \rho_\ee=e^{-c_\ee(\zeta(2+\ee)+P(1))}.
\]
The coefficient is produced by one exceptional action. Deleting that action
from the exactly conditioned Poisson configuration recovers the
unconditioned configuration in total variation. If $nc_\ee\to\infty$, the
remaining actions have a Cauchy-domain cloud on the scale
$b_n=(nc_\ee)^{1/(1+\ee)}$, and the largest action has a reflected, explicitly
normalized $1$-stable fluctuation about an exact truncated-mean center
$D_n$. Consequently, the minimal action cutoff retaining a fixed fraction
$r$ of the coefficient is $D_n+b_nF^{-1}(r)+o(b_n)$. Replacing $D_n$ by
$d_n$ is incorrect at that precision: their difference is asymptotic to
$b_n\log(1/\ee)$. The bounded-intensity regime has instead a discrete
compound-Poisson remainder. We also give a convergent two-coordinate
critical inverse chart, uniform through the removable parameter endpoint,
and explain why its endpoint is a degenerate equation rather than a new
nonzero critical function. All probabilistic estimates needed for the
moving-parameter limits are proved directly.
\end{abstract}
\noindent\textbf{Scope and status.}
This article supplies a model-specific answer to the lower-endpoint part of
Question~4 in the repository's \emph{Critical Hahn Transseries} report.
Cauchy condensation and deletion of a largest summand are established
mechanisms, not discoveries claimed here. The contribution is the explicit
arbitrary-rate triangular theorem, its sparse/dense matching, and its
transseries and action-budget consequences for this family. Publication
priority has not been established. These are conventional mathematical
proofs in an unrefereed research draft; no new Lean verification is claimed.

\clearpage\begingroup\small\setlength{\parskip}{0pt}\tableofcontents\endgroup\clearpage

\section{The precise question and the contribution}
\label{sec:intro}

\subsection{A repository question not answered by the upper endpoint}
The \texttt{ProveIt} transseries project contains a countable-action
critical-inversion model with tail $a_j\asymp j^{-1-\alpha}$,
$1<\alpha<2$. Its report \emph{Beyond Finite-Action Folds: Critical Hahn
Transseries, Stable Sector Asymptotics, and Sharp Action Budgets} asks in
Question~4 for joint limits $\alpha\downarrow1$ and $\alpha\uparrow2$ with
coefficient index $n$ \cite{repo-critical}. The editorial note at the
snapshot inspected here records several results near the upper endpoint,
but explicitly leaves the lower endpoint untreated. The snapshot is
\begin{center}\small\commit.\end{center}
The relevant question occurs at source lines 1502--1507 of the critical
Hahn article; the surrounding editorial note distinguishes the two endpoints.
The separate transseries project and the intake ledger are documented in
\cite{repo-project,repo-index}.

At the upper endpoint the issue is a variance changing from infinite to
finite. At the lower endpoint it is the mean that would cease to be finite
without a simultaneous normalization. The critical normalization tends to
zero. Thus a Gaussian argument with a modified variance cannot simply be
reused, and a fixed-$\alpha$ stable local limit cannot be substituted into a
moving-$\alpha$ problem without new estimates.

We resolve this lower-endpoint problem for an exact power tail with a fixed
finite prefix. The class is broad enough to test universality under local
changes while retaining explicit normalization and exact tail sums. We do
not purport to cover arbitrary slowly varying tails, moving prefixes, or a
simultaneous departure from critical coupling.

\subsection{What belongs to the literature}
The change of variable $t=qe^U$, Lagrange inversion, and conditioning
identities belong to the standard allocation and simply generated tree
framework; Janson gives a comprehensive treatment \cite{janson}. The
principle that one large summand realizes a heavy-tailed deviation, and
that deleting it can restore the original joint law, is developed by
Armend\'ariz and Loulakis \cite{al}. Critical Cauchy condensation, a
separation between the largest component and the remaining extreme-value
scale, and Cauchy fluctuations are already present in the work of
Kortchemski and Richier \cite{kr}. Their setting is a fixed critical
Cauchy-domain offspring law. Here the action law itself changes with $n$,
its exponent approaches one from above, and its total intensity may tend
to infinity, remain bounded, or vanish.

Accordingly, the new result claimed within this article is not the
existence of condensation in general. It is the particular uniform
triangular-array theorem proved below, including exact critical centering,
finite-prefix dependence, and the resulting coefficient and cutoff laws.
We do not invoke a fixed-law condensation theorem as though it were uniform
in our degenerating parameter. The direct estimates are included precisely
to close that logical gap. The special-function input is the classical
polylogarithm expansion recorded in DLMF~25.12.12 \cite{dlmf}.

\subsection{A guide to the results}
The main conclusions are collected in Theorem~\ref{thm:main}. The analytic
chart is constructed in Section~\ref{sec:chart}. Sections~\ref{sec:estimates}
and~\ref{sec:dense} prove the dense local theorem and deletion in total
variation, without any rate restriction on $\ee_n$. Section~\ref{sec:sparse}
handles bounded intensity, including intensity tending to zero. The
Cauchy cloud, action cutoffs, and coefficient matching are established in
Sections~\ref{sec:cloud}--\ref{sec:matching}. Exact algebra checks and
floating-point diagnostics are separated in Section~\ref{sec:computations}.

The strongest reusable ingredient is the deletion theorem. It transfers
\emph{all bounded observables of the remaining configuration}, not just its
largest member or its total weight. The distinction between the leading
condensate scale $d_n$ and the fluctuation center $D_n$ is especially
important for computation: asymptotically equivalent centers need not be
interchangeable in a cutoff quantile.

\section{Model, notation, and exact identities}
\label{sec:model}

\subsection{The finite-prefix class}
Fix a real polynomial
\begin{equation}\label{eq:prefix}
 P(t)=\sum_{j=1}^{J}p_jt^j,
 \qquad \kappa=P'(1)=\sum_{j=1}^{J}jp_j.
\end{equation}
Set $p_j=0$ for $j>J$. Assume that for all sufficiently small real
$\ee>0$,
\begin{equation}\label{eq:positivity}
 j^{-2-\ee}+p_j\ge0\quad(j\ge1),\qquad 1+p_1>0.
\end{equation}
Negative $p_j$ are allowed; it is the resulting weights that must be
nonnegative. The polynomial $P$ is fixed throughout every joint limit.
Write
\begin{align}
 \alpha&=1+\ee,&
 A_\ee(t)&=\Li_{2+\ee}(t)+P(t),\label{eq:A}\\
 c_\ee&=\frac1{\zeta(1+\ee)+\kappa},&
 m_\ee&=c_\ee A_\ee(1),&
 \rho_\ee&=e^{-m_\ee}.\label{eq:c}
\end{align}
For sufficiently small $\ee$ we have $c_\ee>0$, and
\begin{equation}\label{eq:criticalmean}
 c_\ee A_\ee'(1)=1,\qquad
 c_\ee=\ee-(\gamma+\kappa)\ee^2+O(\ee^3).
\end{equation}
Here $\gamma$ denotes Euler's constant.

Define $U_\ee\in q\R[[q]]$ by
\begin{equation}\label{eq:feedback}
 U_\ee(q)=c_\ee A_\ee(qe^{U_\ee(q)}),\qquad
 u_n(\ee)=[q^n]U_\ee(q).
\end{equation}
Each coefficient is determined recursively by lower coefficients. The
analytic implicit function theorem at $(q,U)=(0,0)$ also gives a
convergent solution near the origin. Positivity and the positive first
weight imply $u_n(\ee)>0$ for every $n\ge1$.

For an integer $M\ge0$, let $A_{\ee,M}$ retain only the terms of $A_\ee$
with $j\le M$, and define $U_{\ee,M}$ and $u_{n,M}$ using the \emph{same}
$c_\ee$. Truncation does not renormalize the coupling. This convention is
essential to the cutoff identity below.

\subsection{Poisson configurations}
For each pair $(n,\ee)$, put
\begin{equation}\label{eq:scales}
 \lambda=nc_\ee,\qquad b=\lambda^{1/\alpha},\qquad
 d=(\lambda/\ee)^{1/\alpha}.
\end{equation}
Let the independent variables $N_j$ have laws
\begin{equation}\label{eq:poisson}
 N_j\sim\Pois(\nu_j),\qquad
 \nu_j=\lambda(j^{-1-\alpha}+p_j),\qquad
 K=\sum_{j\ge1}jN_j.
\end{equation}
We extend $\nu_j=0$ for integer $j\le0$. A restricted power multiplied
by an event indicator is evaluated on that event and defined to be zero
outside it. Since $\sum_j\nu_j<\infty$, the count vector has finite support
almost surely. Its total weight satisfies $\E K=n$. Conditioned on $K=n$, list the
actions, with multiplicity, in nonincreasing order:
\[
 J_1\ge J_2\ge\cdots\ge0.
\]
Zeros pad a finite list. Let $\mathcal D(N)$ be the count vector after
removing one occurrence of a largest action; the choice between equal
largest occurrences does not change the count vector.

We use $\tv(\mu,\nu)=\sup_E|\mu(E)-\nu(E)|$ for total variation distance.
The space of finitely supported count vectors on the positive integers is
countable, so no additional measurable-space construction is required.
All conditional assertions about $J_i$ refer to $\Prob(\,\cdot\mid K=n)$.

\begin{proposition}[Exact coefficient and cutoff identities]\label{prop:exact}
With $t=qe^{U_\ee(q)}$,
\begin{equation}\label{eq:parametric}
 U_\ee=c_\ee A_\ee(t),\qquad q=t e^{-c_\ee A_\ee(t)},
 \qquad
 u_n(\ee)=\frac1n[t^n]e^{nc_\ee A_\ee(t)}.
\end{equation}
Consequently,
\begin{equation}\label{eq:coefficientprob}
 u_n(\ee)=\frac{\rho_\ee^{-n}}n\Prob(K=n),\qquad
 \frac{u_{n,M}(\ee)}{u_n(\ee)}
 =\Prob(J_1\le M\mid K=n).
\end{equation}
The ratio is nondecreasing in $M$ and is exactly one for $M\ge n$.
\end{proposition}
\begin{proof}
The first two identities are direct substitutions. Lagrange inversion for
$t=q\exp(c_\ee A_\ee(t))$ gives
\[
 [q^n]c_\ee A_\ee(t(q))
 =\frac1n[t^{n-1}]c_\ee A_\ee'(t)e^{nc_\ee A_\ee(t)}
 =\frac1n[t^n]e^{nc_\ee A_\ee(t)}.
\]
For the last equality, differentiate the exponential and compare its
coefficient of $t^{n-1}$. The probability generating function of $K$ is
$\exp(nc_\ee(A_\ee(t)-A_\ee(1)))$, proving the first probability identity.
Forbidding all actions above $M$ multiplies the truncated Poisson
probability by $\exp(-\sum_{j>M}\nu_j)$. Thus its generating function is
$\exp(nc_\ee A_{\ee,M}(t)-nc_\ee A_\ee(1))$, proving the second identity.
Monotonicity and exactness for $M\ge n$ follow from the event description.
\end{proof}

\section{Main theorem: arbitrary-rate approach to the lower endpoint}
\label{sec:main}
Throughout a limit means $n\to\infty$ through integers and
$\ee=\ee_n\to0^+$, unless a different limit is explicitly specified.
Subscripts on $b,d,\lambda$ are often suppressed. No monotonicity of
$\ee_n$ is required.

Define the real random variable $Z$ by its characteristic function
\begin{align}
 \E e^{itZ}
 &=\exp\left\{\int_0^\infty
       (e^{itx}-1-itx\ind_{x\le1})\frac{\dd x}{x^2}\right\}\label{eq:stablecf}\\
 &=\exp\left\{-\frac\pi2|t|+it(1-\gamma-\log|t|)\right\}
 \quad(t\ne0),\nonumber
\end{align}
with value one at $t=0$. This is a totally right-skewed $1$-stable law
with the displayed location convention; it is not a symmetric Cauchy
law. Write
\begin{equation}\label{eq:F}
 F(x)=\Prob(-Z\le x).
\end{equation}
The function $F$ is continuous and strictly increasing from zero to one,
as proved in Section~\ref{sec:cloud}.

\begin{theorem}[Uniform coefficient, deletion, and cutoff laws]\label{thm:main}
For the fixed-prefix family \eqref{eq:prefix}--\eqref{eq:feedback} the
following statements hold.

\smallskip\noindent\textbf{(a) Every rate of approach.}
For every $\ee_n\to0^+$,
\begin{equation}\label{eq:maincoefficient}
 \boxed{\displaystyle
 u_n(\ee_n)\sim\rho_{\ee_n}^{-n}\frac{\ee_n}{n d_n}.}
\end{equation}
In the conditioned configuration,
\begin{equation}\label{eq:mainTV}
 \tv\bigl(\law(\mathcal D(N)\mid K=n),\law(N)\bigr)\longrightarrow0,
 \qquad \frac{J_1}{d_n}\pto1,\qquad \frac{J_2}{J_1}\pto0.
\end{equation}

\smallskip\noindent\textbf{(b) Diverging intensity.}
Suppose $\lambda_n=nc_{\ee_n}\to\infty$. Put
\begin{equation}\label{eq:D}
 \mu_b=\sum_{j\le b}j\nu_j,\qquad D_n=n-\mu_b.
\end{equation}
For all sufficiently large $n$,
\begin{equation}\label{eq:Dtail}
 D_n=\lambda_n\zeta(1+\ee_n,\lfloor b_n\rfloor+1)
     =\frac{b_n}{\ee_n}+O(1).
\end{equation}
Then
\begin{equation}\label{eq:mainfluct}
 \frac{J_1-D_n}{b_n}\dto-Z,
 \qquad
 \sum_{k\ge2}\delta_{J_k/b_n}\dto\Pi,
\end{equation}
where $\Pi$ is a Poisson point process on $(0,\infty)$ with intensity
$x^{-2}\dd x$, with convergence of point measures away from zero.
The two limits hold jointly, with $Z$ the compensated sum of this same
point process. In particular, for $x\in\R$ and $r\in(0,1)$,
\begin{align}
 \frac{u_{n,\lfloor D_n+b_nx\rfloor}(\ee_n)}{u_n(\ee_n)}
 &\longrightarrow F(x),\label{eq:maincutoff}\\
 M_r(n):=\min\{M\in\N:u_{n,M}/u_n\ge r\}
 &=D_n+b_nF^{-1}(r)+o(b_n).\label{eq:mainbudget}
\end{align}

\smallskip\noindent\textbf{(c) Bounded intensity.}
If $\lambda_n\to\lambda_0\in[0,\infty)$, let independent
$N_j^0\sim\Pois(\lambda_0(j^{-2}+p_j))$ and
$S_{\lambda_0}=\sum_jjN_j^0$. Then
\begin{equation}\label{eq:mainsparse}
 n-J_1\dto S_{\lambda_0},\qquad
 \Prob(K=n)\sim\lambda_n n^{-2-\ee_n}.
\end{equation}
The first convergence holds in total variation as a convergence of
integer-valued laws. If $\lambda_0=0$, $J_1=n$ with probability tending to
one.
\end{theorem}

\begin{remark}[Meaning of arbitrary-rate uniformity]\label{rem:uniform}
Part (a) holds along every pair sequence $n\to\infty$, $\ee\to0^+$.
Equivalently, for every $\eta>0$ there are $N$ and $\ee_0>0$ such that the
relative error in \eqref{eq:maincoefficient} is below $\eta$ whenever
$n\ge N$ and $0<\ee\le\ee_0$. This equivalence is the elementary
sequential characterization of a two-variable limit. It gives existence
of uniform thresholds, not effective numerical values of those thresholds.
\end{remark}

\begin{remark}[What is and is not universal]
In (b) the limiting standardized cloud is independent of the finite
prefix. The normalizers $c_\ee,\rho_\ee,b_n,D_n$ still depend on it.
In (c) the limiting discrete cloud retains all finite-prefix weights.
Neither statement says that changing a prefix leaves raw coefficients
asymptotically unchanged: the factor $\rho_\ee^{-n}$ can magnify a small
change exponentially.
\end{remark}

\section{A convergent critical inverse chart through the degenerating endpoint}
\label{sec:chart}

\subsection{Exact phase and removable singularities}
Take $t=e^{-x}$, with a fixed branch of $\log x$ near the positive real
axis, and let $h=\log(\rho_\ee/q)$. The exact parametrization gives
\begin{equation}\label{eq:phaseexact}
 h=x+c_\ee(A_\ee(e^{-x})-A_\ee(1)),\qquad
 U_\ee=m_\ee+h-x.
\end{equation}
The polylogarithm expansion \cite[Eq.~25.12.12]{dlmf} yields
\begin{equation}\label{eq:phase}
 h=B_\ee x^{1+\ee}+c_\ee x^2 R_\ee(x),\qquad
 B_\ee=c_\ee\Gamma(-1-\ee),
\end{equation}
where
\begin{equation}\label{eq:R}
 R_\ee(x)=\sum_{k\ge2}\frac{(-1)^k}{k!}
 \left(\zeta(2+\ee-k)+\sum_{j=1}^Jp_jj^k\right)x^{k-2}.
\end{equation}
The radius in $x$ is at least $2\pi$ before choosing a smaller common
polydisc. In particular the linear term in \eqref{eq:phaseexact} vanishes
\emph{exactly}, by criticality.

\begin{lemma}[Uniform analytic coefficients]\label{lem:analytic}
There is $e_0>0$ such that $c_\ee$, $B_\ee$, and $R_\ee(x)$ extend
holomorphically in $\ee$ to $|\ee|<e_0$, with $R$ jointly holomorphic for
$|x|<r_0$ for some $r_0>0$. Moreover
\begin{equation}\label{eq:Bexp}
 c_0=0,\qquad B_0=1,\qquad
 B_\ee=1-(1+\kappa)\ee+O(\ee^2).
\end{equation}
After shrinking $e_0$, $B_\ee$ has no zero there.
\end{lemma}
\begin{proof}
Write $\zeta(1+\ee)=\ee^{-1}+\gamma+O(\ee)$ and
$\Gamma(-1-\ee)=\ee^{-1}+\gamma-1+O(\ee)$.
The first reciprocal has a simple zero at zero and cancels the latter
pole. Multiplication proves \eqref{eq:Bexp}. The coefficients of
\eqref{eq:R} have no pole near $\ee=0$. The convergent polylogarithm
expansion, or its locally uniformly convergent power-series remainder,
shows joint holomorphy in a smaller polydisc. Nonvanishing follows from
$B_0=1$ and continuity.
\end{proof}

\subsection{Uniform two-coordinate inversion}
Set $\delta=1-\ee=2-\alpha$ and
\begin{equation}\label{eq:chartcoordinates}
 y=(h/B_\ee)^{1/\alpha},\qquad
 z=\frac{c_\ee}{B_\ee}y^\delta
   =\frac{y^\delta}{\Gamma(-1-\ee)}.
\end{equation}
Branches are taken consistently with positive $h,x,y$ for positive real
arguments. The variables $y,z$ in the next theorem are first independent
complex coordinates; only afterwards are they specialized as above.

\begin{theorem}[Uniform convergent inverse chart]\label{thm:chart}
There are positive radii $r,s,e_1$ and a unique holomorphic function
$V(y,z,\ee)$ near $\{|y|<r,|z|<s,|\ee|<e_1\}$, with
$|V-1|<1/4$, satisfying
\begin{equation}\label{eq:Veq}
 V^\alpha\bigl(1+zV^\delta R_\ee(yV)\bigr)=1,
 \qquad V(y,0,\ee)=1.
\end{equation}
In these coordinates the local inverse is
\begin{equation}\label{eq:chart}
 x=yV(y,z,\ee),\qquad
 U_\ee=m_\ee+h-yV(y,z,\ee).
\end{equation}
The double power series of $V$ converges absolutely on a common smaller
polydisc. Thus, for each sufficiently small real $\ee>0$, the specialized
$x$ is a convergent Hahn series on the positive grid
\begin{equation}\label{eq:haangrid}
 \{1+k+m(1-\ee):k,m\in\N\}.
\end{equation}
The coefficient functions extend holomorphically through $\ee=0$.
\end{theorem}
\begin{proof}
At $V=1,z=0$ the left side of \eqref{eq:Veq} minus one vanishes, and its
$V$ derivative is $\alpha$, bounded away from zero for small complex
$\ee$. The holomorphic implicit function theorem supplies the chart,
uniformly after restricting $|y|,|\ee|$ to smaller closed discs. Equivalently,
consider
\[
 T(V)=\bigl(1+zV^\delta R_\ee(yV)\bigr)^{-1/\alpha}.
\]
On $|V-1|\le1/4$, small common radii bound $R$, its derivative, and the
chosen powers. For sufficiently small $s$, $T$ maps this closed disc into
its interior and has derivative of modulus at most $1/2$. This proves
uniqueness in the stated disc and gives a convergent iteration there.
Substituting $x=yV$ into \eqref{eq:phase} gives exactly \eqref{eq:Veq}.
Ordinary Cauchy estimates on a smaller polydisc prove absolute convergence
of the double power series. Specialization of $z$ then produces the grid
\eqref{eq:haangrid}. Its generators are strictly positive, so only finitely
many pairs contribute below any fixed exponent bound. Coinciding
exponents are simply collected. Holomorphic dependence on $\ee$ was part
of the implicit-function construction.
\end{proof}

\begin{proposition}[First chart coefficients and a geometric error bound]
\label{prop:chartcoeff}
For fixed $\ee$, write $R=R_\ee(y)$ and $R'=R_\ee'(y)$. Then
\begin{equation}\label{eq:Vcoeff}
 V=1-\frac{R}{\alpha}z+
 \frac{(5-\alpha)R^2+2yRR'}{2\alpha^2}z^2+O(z^3),
\end{equation}
uniformly on smaller $(y,\ee)$ polydiscs. Suppose $|V|\le H$ on
$|y|\le r$, $|z|\le s$ at a fixed admissible parameter. If $V_{K,M}$
retains $0\le k\le K$, $0\le m\le M$ in the double Taylor series, then
for $a=|y|/r<1$ and $v=|z|/s<1$,
\begin{equation}\label{eq:chartbound}
 |V-V_{K,M}|
 \le \frac{H(a^{K+1}+v^{M+1})}{(1-a)(1-v)}.
\end{equation}
The bound is uniform whenever the same $H,r,s$ work for the parameters.
\end{proposition}
\begin{proof}
Using $\alpha+\delta=2$, rewrite \eqref{eq:Veq} as
$V^\alpha+zV^2R_\ee(yV)=1$. Substitute $V=1+az+bz^2+O(z^3)$.
The coefficients of $z$ and $z^2$ respectively give
\[
 \alpha a+R=0,\qquad
 \alpha b+\frac{\alpha(\alpha-1)}2a^2+a(2R+yR')=0,
\]
which prove \eqref{eq:Vcoeff}. Cauchy's estimate gives
$|[y^kz^m]V|\le Hr^{-k}s^{-m}$. Sum the two geometric tails outside the
retained rectangle; counting their intersection twice only increases the
upper bound.
\end{proof}

\subsection{The endpoint is degenerate, not a hidden nonzero solution}
At $\ee=0$ we have $c_0=0$, $B_0=1$, and the specialized $z$ is zero.
Hence $h=x$, $m_0=0$, and \eqref{eq:chart} gives $U_0=0$. This does not
contradict the nontrivial coefficient theorem: $n$ grows while the parameter
approaches its degenerate endpoint.

\begin{proposition}[Fixed-domain limit]\label{prop:fixedq}
On each closed disc $|q|\le r<1$,
\begin{equation}\label{eq:fixedq}
 \frac{U_\ee(q)}{c_\ee}\longrightarrow\Li_2(q)+P(q)
\end{equation}
uniformly as $\ee\downarrow0$.
\end{proposition}
\begin{proof}
Choose $r<r'<1$. For $|q|\le r$ and $|U|$ small enough,
$|qe^U|\le r'$. The functions $A_\ee$ and their derivatives are uniformly
bounded on this latter disc and converge there uniformly to $A_0$.
For small $c_\ee$, the map $U\mapsto c_\ee A_\ee(qe^U)$ is a uniform
contraction on a disc of radius $C c_\ee$. Consequently $U_\ee=O(c_\ee)$
uniformly. Divide the equation by $c_\ee$ and pass to the uniform limit.
\end{proof}

\section{Uniform estimates for the Poisson array}
\label{sec:estimates}
This section supplies the local estimates that cannot be obtained merely
by substituting a moving exponent into a fixed-parameter theorem. Constants
may depend on the fixed prefix, but not on $n$ or on sufficiently small
$\ee$. Eventually we restrict $1<\alpha\le5/4$.

\subsection{Tail sums and truncated variance}
For $M>J$, define
\begin{equation}\label{eq:truncated}
 K_M=\sum_{j\le M}jN_j,\quad
 \Lambda_M=\sum_{j>M}\nu_j,\quad
 \Delta_M=n-\E K_M,\quad V_M=\Var K_M.
\end{equation}
The integral test gives, uniformly in the indicated range,
\begin{align}
 \Lambda_M&=\frac{\lambda M^{-\alpha}}{\alpha}
             +O(\lambda M^{-1-\alpha}),\label{eq:tailnumber}\\
 \Delta_M&=\frac{\lambda M^{-\ee}}{\ee}
             +O(\lambda M^{-1-\ee}),\label{eq:tailmean}\\
 V_M&\le C\lambda(M^{2-\alpha}+1),\qquad
 \sup_{j>M}\nu_j\le C\lambda M^{-1-\alpha}.\label{eq:variance}
\end{align}
The mean identity uses $\E K=n$. Even when some $p_j$ are negative, the
variance is bounded above using $\sum_{j\le J}|p_j|j^2$ plus the pure tail.
If $b\to\infty$, \eqref{eq:tailmean} at $M=\lfloor b\rfloor$ gives
\begin{equation}\label{eq:Dproof}
 D=\lambda\sum_{j>\lfloor b\rfloor}j^{-1-\ee}
   =\frac{b}{\ee}+O(1),
\end{equation}
because $\lambda=b^{1+\ee}$. The floor contributes only $O(1)$ as well.

\begin{lemma}[Convolution bound for two or more large actions]
\label{lem:twobig}
Let $\nu^>$ be the measure on integers with mass $\nu_j$ at $j>M$,
let $\Lambda=\|\nu^>\|_1$, and let $H_M=\sum_{j>M}jN_j$. Then
\begin{equation}\label{eq:twobig}
 \sup_k\Prob(K_M+H_M=k,\ \text{at least two actions exceed }M)
 \le \|\nu^>\|_\infty\Lambda e^\Lambda.
\end{equation}
\end{lemma}
\begin{proof}
The subprobability measure for exactly $l$ large actions is
$e^{-\Lambda}(\nu^>)^{*l}/l!$. For $l\ge2$,
$\|(\nu^>)^{*l}\|_\infty\le\|\nu^>\|_\infty\Lambda^{l-1}$.
Convolution with the probability law of $K_M$ cannot increase the supremum
norm. Sum over $l\ge2$ and use
$\sum_{l\ge2}\Lambda^{l-1}/l!\le\Lambda e^\Lambda$.
\end{proof}

\subsection{A local bound preserved by positive exponential tilting}
\begin{lemma}[Uniform lattice concentration]\label{lem:concentration}
Assume $b\to\infty$ and $M/b\to\infty$. For every $t\ge0$, let
$\widetilde K_M$ be compound Poisson with rates
$\widetilde\nu_j=\nu_je^{tj}$, $j\le M$. There is a constant $C$ such that
\begin{equation}\label{eq:lattice}
 \sup_{k\in\mathbb Z}\Prob(\widetilde K_M=k)\le\frac Cb
\end{equation}
for all sufficiently large $n$, uniformly in $t$.
\end{lemma}
\begin{proof}
Fourier inversion bounds the left side by
\[
 \frac1{2\pi}\int_{-\pi}^{\pi}
 \exp\left(-\sum_{j\le M}\nu_je^{tj}(1-\cos j\theta)\right)\dd\theta.
\]
Since $t\ge0$, it suffices to bound the untilted integral. The interval
$|\theta|\le b^{-1}$ contributes at most $2/b$. Choose a small fixed
$\theta_0>0$ so that, for $b^{-1}\le|\theta|\le\theta_0$, the integers
$j$ in $[1/(2|\theta|),1/|\theta|]$ include at least a constant multiple
of $1/|\theta|$ indices above $J$. They are at most $b$ and hence below
$M$. On this range, $1-\cos(j\theta)$ is bounded below by a positive
constant. The exact tail gives
\[
 \sum_{j\le M}\nu_j(1-\cos j\theta)
 \ge c\lambda|\theta|^\alpha=c(b|\theta|)^\alpha.
\]
The integral of this upper envelope is $O(1/b)$ uniformly for
$1\le\alpha\le5/4$.

On $\theta_0\le|\theta|\le\pi$, retain two consecutive fixed indices
larger than $J$. Their trigonometric terms have no common zero in this
compact set, and both rates are at least a positive constant times
$\lambda$. The remaining integral is $O(e^{-c\lambda})=O(1/b)$.
All estimates survive the positive tilt. This proves the claim.
\end{proof}

\begin{lemma}[Tilt bound at a prescribed lattice point]\label{lem:tilt}
Under the hypotheses of Lemma~\ref{lem:concentration}, for $t\ge0$,
\begin{equation}\label{eq:tiltbound}
 \Prob(K_M=n)\le\frac Cb
 \exp\left(-t\Delta_M+\frac{V_M}{M^2}e^{tM}\right).
\end{equation}
Without the concentration hypothesis, the corresponding exponential
factor bounds the tail probability $\Prob(K_M\ge n)$.
\end{lemma}
\begin{proof}
The exact exponential-tilt identity at $n$ contributes the factor
$\exp(-tn+\sum_{j\le M}\nu_j(e^{tj}-1))$ multiplying the tilted local
probability. For $0\le j\le M$, the power series of the exponential gives
\[
 e^{tj}-1-tj\le\frac{j^2}{M^2}(e^{tM}-1-tM)
               \le\frac{j^2}{M^2}e^{tM}.
\]
Sum this inequality and use Lemma~\ref{lem:concentration}. The tail bound
is the usual exponential Markov inequality with the same calculation.
\end{proof}

\section{Dense condensation and deletion in total variation}
\label{sec:dense}
Assume throughout this section that $\lambda\to\infty$, equivalently
$b\to\infty$. The proof does not require $\ee\log n$ to be bounded.
Define the positive real tail weight
\begin{equation}\label{eq:wD}
 w_D=\lambda D^{-1-\alpha}.
\end{equation}
This is evaluated at the real number $D$; replacing it by the weight at
$\lfloor D\rfloor$ has relative error tending to zero.

\begin{theorem}[Dense local one-action theorem]\label{thm:dense}
For every sequence with $\lambda\to\infty$ and $\ee\to0^+$,
\begin{equation}\label{eq:denselocal}
 \Prob(K=n)\sim w_D.
\end{equation}
With conditional probability tending to one, exactly one action exceeds
$\eta D$, for any one fixed sufficiently small $\eta>0$. That action,
divided by $D$, tends to one in probability. Moreover the deletion
conclusion \eqref{eq:mainTV} holds.
\end{theorem}
\begin{proof}
Fix $\eta=1/64$ and take $M=\lfloor\eta D\rfloor$. Since
$D=b/\ee+O(1)$,
\begin{equation}\label{eq:densebounds}
 \frac Mb\longrightarrow\infty,\qquad
 \Lambda_M=O(\ee^\alpha)=o(1),\qquad
 \sup_{j>M}\nu_j=O(w_D),\qquad
 \frac{V_M}{D^2}=O(\ee^\alpha).
\end{equation}
Furthermore, using \eqref{eq:tailmean},
\begin{equation}\label{eq:Deltaratio}
 \frac{\Delta_M}{D}
 =\left(\frac bM\right)^\ee(1+o(1))
 =\left(\frac\ee\eta\right)^\ee(1+o(1))\longrightarrow1.
\end{equation}
Chebyshev's inequality therefore gives
\begin{equation}\label{eq:typicalsmall}
 \frac{n-K_M}{D}\pto1.
\end{equation}
This statement concerns the unconditioned small-action configuration.

\smallskip\noindent\emph{No large action.}
Set $tM=\tfrac12\log(1/\ee)$. The positive remainder in the exponent
of \eqref{eq:tiltbound} is bounded by
\[
 C\lambda(M^{2-\alpha}+1)M^{-2}\ee^{-1/2}
 =O(\ee^{\alpha-1/2})+o(1)=o(1).
\]
In the second term one may use
$\lambda/M^2=b^{\alpha-2}O(\ee^2)$ and $\alpha\le5/4$.
Meanwhile $t\Delta_M=(1/(2\eta)+o(1))\log(1/\ee)$.
Thus for all sufficiently large $n$,
\begin{equation}\label{eq:nobig}
 \Prob(K_M=n)\le\frac Cb\ee^{16}=o(w_D),
\end{equation}
because $w_D\sim\ee^{1+\alpha}/b$. The joint probability with no large
action is still smaller. The factor $1/b$ in this estimate is crucial;
a bound only on the upper tail would not suffice for every possible rate
of growth of $b$.

\smallskip\noindent\emph{At least two large actions.}
By Lemma~\ref{lem:twobig} and \eqref{eq:densebounds}, the local probability
of $K=n$ with at least two actions above $M$ is
$O(w_D\Lambda_M)=o(w_D)$.

\smallskip\noindent\emph{Exactly one large action.}
Its contribution is exactly
\begin{equation}\label{eq:onebigexact}
 e^{-\Lambda_M}\E\left[
   \nu_{n-K_M}\ind_{n-K_M>M}\right].
\end{equation}
For large $n$ the prefix lies below $M$. Dividing the integrand by $w_D$
gives a uniformly bounded nonnegative variable:
\[
 \frac{\nu_{n-K_M}}{w_D}\ind_{n-K_M>M}
 =\left(\frac{n-K_M}{D}\right)^{-1-\alpha}
       \ind_{n-K_M>M}\le C_\eta.
\]
By \eqref{eq:typicalsmall} it converges to one in probability and hence in
$L^1$. Since $\Lambda_M\to0$, \eqref{eq:onebigexact} is asymptotic to
$w_D$. The three disjoint cases prove \eqref{eq:denselocal} and show that
exactly one large action accounts for asymptotically all the conditioned
probability. Its size divided by $D$ tends to one by the same bounded
weight argument.

\smallskip\noindent\emph{Deletion.}
On the event just considered, the conditional law of the small count
vector has density
\begin{equation}\label{eq:RN}
 \frac{\nu_{n-K_M}\ind_{n-K_M>M}}
      {\E[\nu_{n-K_M}\ind_{n-K_M>M}]}
\end{equation}
relative to its original Poisson law. This density converges to one in
$L^1$, so the two laws converge in total variation. The original full
configuration differs from the original small configuration with
probability at most $1-e^{-\Lambda_M}\to0$, under their natural coupling.
The exceptional conditioned events have probability tending to zero.
Finally, on the exactly-one-large-action event the deleted action is the
largest action. The triangle inequality for total variation proves the
claim.
\end{proof}

\begin{remark}[Why critical mean does not prevent condensation]
Every row has $\E K=n$, but its typical value is $\mu_b+O_{\Prob}(b)$,
not $n$. Their difference is $D\sim b/\ee\gg b$. Very rare actions beyond
the typical cloud carry the missing expectation. Conditioning on $K=n$
selects one such action. There is no contradiction between deletion in
total variation and a changed expectation, since total variation controls
bounded observables and these weights are not uniformly integrable.
The fluctuation assertion used in this explanation is proved next in
Section~\ref{sec:cloud}; it is not needed in the local proof above.
\end{remark}

\section{Bounded intensity and the discrete remainder}
\label{sec:sparse}
The regime $\lambda=nc_\ee=O(1)$ is not a small variation on a stable
local limit: $b$ need not diverge and can tend to zero. In particular a
Fourier bound proportional to $1/b$ would be useless here. The following
argument also keeps a factor of $\lambda$ when $\lambda\to0$, which is
necessary for relative asymptotics.

\begin{theorem}[Sparse local theorem]\label{thm:sparse}
Uniformly along sequences with $\lambda=O(1)$ and $\ee\to0^+$,
\begin{equation}\label{eq:sparselocal}
 \Prob(K=n)\sim\lambda n^{-1-\alpha}.
\end{equation}
Deletion of one largest action again gives the original Poisson
configuration in total variation. If $\lambda\to\lambda_0$, the conclusion
\eqref{eq:mainsparse} follows.
\end{theorem}
\begin{proof}
Since $c_\ee/\ee\to1$, bounded $\lambda$ implies $\ee=O(1/n)$ and hence
$\ee\log n\to0$. Set $M=\lfloor\eta n\rfloor$ with $\eta=1/64$.
Then
\begin{equation}\label{eq:sparseestimates}
 \E K_M\le C\lambda\log n,\quad
 V_M\le C\lambda n,\quad
 \Lambda_M=O(\lambda/n),\quad
 \sup_{j>M}\nu_j=O(\lambda n^{-1-\alpha}).
\end{equation}
In particular $K_M/n\pto0$. Lemma~\ref{lem:twobig} makes the local
contribution of at least two large actions negligible relative to
$\lambda n^{-1-\alpha}$. Exactly one large action contributes
\[
 e^{-\Lambda_M}\E[\nu_{n-K_M}\ind_{n-K_M>M}]
 \sim\lambda n^{-1-\alpha},
\]
by the same bounded-weight argument used in the dense case.

It remains to estimate zero large actions \emph{relatively}, even when
$\lambda$ is arbitrarily small. Let $tM=\tfrac12\log n$ and put
$L=\sum_{j\le M}\nu_j$, $L_t=\sum_{j\le M}\nu_je^{tj}$.
The exponential estimate used before now shows
\[
 L_t-L
 \le t\E K_M+\frac{V_M}{M^2}e^{tM}
 \le C\lambda\frac{\log^2n}{n}+C\lambda n^{-1/2}=o(1).
\]
For $n>0$, the zero configuration cannot contribute. Therefore
\begin{align*}
 \Prob(K_M=n)
 &\le e^{-tn}\E[e^{tK_M};K_M>0]\\
 &=e^{-tn-L}(e^{L_t}-1)
 \le e^{-tn+L_t-L}L_t
 \le C\lambda n^{-1/(2\eta)+1/2}.
\end{align*}
Here $L_t\le e^{tM}L\le C\lambda\sqrt n$.
The final bound is $o(\lambda n^{-1-\alpha})$; we used
$\ee\log n\to0$ to compare $n^{-1-\alpha}$ with $n^{-2}$.
This proves the local equivalent without dividing by an uncontrolled
small probability.

The Radon--Nikodym density for the small configuration, conditioned on
exactly one large action and total weight $n$, is again
\eqref{eq:RN}. Its normalization is now $\lambda n^{-1-\alpha}$.
The density converges to one in $L^1$, and $\Lambda_M\to0$, proving
convergence after deletion in total variation.

Finally, if $\lambda\to\lambda_0$, dominated convergence gives
\[
 \sum_{j\ge1}\left|
 \lambda(j^{-2-\ee}+p_j)-\lambda_0(j^{-2}+p_j)\right|\longrightarrow0.
\]
Independent Poisson configurations whose rate measures differ by this
$\ell^1$ amount can be coupled to agree except on an event of probability
at most that amount, by using shared rates and independent excess rates.
Thus the original full configuration converges in total variation to
$(N_j^0)$. Deletion and the map $N\mapsto\sum jN_j$ preserve the resulting
upper bound on total variation, and the total weight after deletion is
exactly $n-J_1$. This proves the remaining assertions.
\end{proof}

\begin{remark}[Finite almost surely, infinite mean]
For $\lambda_0>0$, the limiting sparse cloud has finite total number of
actions almost surely because $\sum_j(j^{-2}+p_j)<\infty$. Its total
weight is nevertheless unbounded and has infinite mean, since
$\sum_j j(j^{-2}+p_j)=\infty$. There is no assertion of moment convergence
in \eqref{eq:mainsparse}. When $\lambda_0=0$, its configuration is empty
almost surely.
\end{remark}

\begin{corollary}[Discrete sparse cutoff law]\label{cor:sparsecutoff}
If $\lambda\to\lambda_0\in[0,\infty)$, then for every fixed integer
$k\ge0$,
\begin{equation}\label{eq:sparsecutoff}
 \frac{u_{n,n-k}(\ee)}{u_n(\ee)}
 \longrightarrow\Prob(S_{\lambda_0}\ge k).
\end{equation}
For $\lambda_0=0$ the limits are one for $k=0$ and zero for each $k\ge1$.
\end{corollary}
\begin{proof}
The event $J_1\le n-k$ is the event $n-J_1\ge k$. Apply the integer-law
total variation convergence from Theorem~\ref{thm:sparse}.
\end{proof}

\section{The Cauchy cloud and conditional extremes}
\label{sec:cloud}
Return to the dense regime $\lambda\to\infty$. The deletion theorem
reduces the conditional limit laws to an unconditioned Poisson limit,
which we now prove with the required moving exponent.

\subsection{Joint convergence of points and compensated sums}
\begin{theorem}[Unconditioned triangular Cauchy limit]\label{thm:cloud}
If $b\to\infty$ and $\ee\to0^+$, then
\begin{equation}\label{eq:cloud}
 \left(\sum_{j\ge1}N_j\delta_{j/b},\frac{K-\mu_b}{b}\right)
 \dto(\Pi,Z),
\end{equation}
where $\Pi$ has intensity $x^{-2}\dd x$ and $Z$ is its compensated sum:
\begin{equation}\label{eq:compensated}
 Z=\lim_{a\downarrow0}
 \left(\sum_{x\in\Pi,\ x>a}x-\log(1/a)\right).
\end{equation}
The small-point compensation converges in $L^2$; the finitely many points
above one are left uncompensated. The point process and $Z$ are not
independent.
\end{theorem}
\begin{proof}
For a continuous function $f$ of compact support in $(0,\infty)$, eventually
all indices in its support lie above the prefix. Since $\lambda=b^\alpha$,
\begin{equation}\label{eq:Riemann}
 \sum_j\nu_jf(j/b)
 =\frac1b\sum_j (j/b)^{-1-\alpha}f(j/b)
 \longrightarrow\int_0^\infty f(x)x^{-2}\dd x.
\end{equation}
This is a Riemann sum, uniformly on compact intervals separated from
zero, and $\alpha\to1$. The Poisson Laplace functional
$\exp(\sum_j\nu_j(e^{-f(j/b)}-1))$ proves convergence of point measures
on those intervals. In addition,
\begin{equation}\label{eq:largerpoints}
 \Prob(\text{some action exceeds }Rb)
 \le\sum_{j>Rb}\nu_j\le C R^{-\alpha}\le C/R\quad(R\ge1)
\end{equation}
for all sufficiently large $n$. This controls the passage to unbounded
intervals at the upper end.

Fix $0<a<1$. The contribution from points above $a$ to the centered sum is
\[
 \sum_{j>ab}\frac jbN_j-\sum_{ab<j\le b}\frac jb\nu_j.
\]
Together with the point process, it converges by \eqref{eq:Riemann} and
\eqref{eq:largerpoints} to
$\sum_{x\in\Pi,x>a}x-\int_a^1x^{-1}\dd x$.
The error from the centered points below $a$ has variance
\begin{align*}
 \frac1{b^2}\sum_{j\le ab}j^2\nu_j
 &\le C a^{2-\alpha}+C\lambda/b^2\\
 &\le C a^{3/4}+o(1),
\end{align*}
where the second term bounds the fixed prefix and tends to zero because
$\lambda/b^2=b^{\alpha-2}$. This variance bound tends to zero after first
letting $n\to\infty$ and then $a\downarrow0$.
The analogous limiting variance is
$\int_0^a x^2x^{-2}\dd x=a$; hence the compensated small-point sum
converges in $L^2$. Truncation at $a$ followed by these variance estimates
proves the joint convergence and the construction \eqref{eq:compensated}.
Its characteristic function is precisely \eqref{eq:stablecf}.
\end{proof}

\begin{proposition}[Normalization and regularity of the stable law]
\label{prop:stable}
The integral exponent in \eqref{eq:stablecf} equals
$-\pi|t|/2+it(1-\gamma-\log|t|)$. The distribution of $Z$ has a continuous
real-analytic density, and its distribution function, and hence $F$, is
strictly increasing.
\end{proposition}
\begin{proof}
For real $s>0$, differentiation under convergent truncated integrals gives
\[
 \frac{\mathrm d}{\mathrm ds}
 \int_0^\infty(e^{-sx}-1+sx\ind_{x\le1})\frac{\dd x}{x^2}
 =\log s+\gamma.
\]
For example, the constant at $s=1$ is the familiar integral
$\int_0^1(1-e^{-x})\dd x/x-\int_1^\infty e^{-x}\dd x/x=\gamma$;
the derivative of the right side with respect to $s$ is $1/s$.
Integration from zero gives $s\log s+(\gamma-1)s$.
Analytic continuation to $s=-it$ from $\Rea s>0$, or direct dominated
limits in the compensated integral, gives the stated Fourier exponent.
The continuous extension at zero follows from $t\log|t|\to0$.

The characteristic function has modulus $e^{-\pi|t|/2}$. Fourier
inversion therefore supplies a continuous density that extends
holomorphically to a strip of width less than $\pi/2$ about the real axis.
It is nonnegative and has integral one. A real-analytic function that
vanishes throughout an interval must vanish identically; consequently this
density has positive integral on every nonempty interval. The distribution
function is continuous and strictly increasing, with limits zero and one.
\end{proof}

\subsection{Transfer through deletion}
\begin{theorem}[Largest action and the remaining extremes]\label{thm:extremes}
In the dense regime, under exact conditioning $K=n$,
\begin{equation}\label{eq:extremes}
 \left(\sum_{k\ge2}\delta_{J_k/b},\frac{J_1-D}{b}\right)
 \dto(\Pi,-Z).
\end{equation}
For every $x>0$,
\begin{equation}\label{eq:secondlargest}
 \Prob(J_2/b\le x\mid K=n)\longrightarrow e^{-1/x}.
\end{equation}
More generally, for each integer $k\ge1$,
\begin{equation}\label{eq:kthlargest}
 \Prob(J_{k+1}/b\le x\mid K=n)
 \longrightarrow e^{-1/x}\sum_{l=0}^{k-1}\frac{x^{-l}}{l!}.
\end{equation}
\end{theorem}
\begin{proof}
Total variation in Theorem~\ref{thm:dense} permits every measurable
$n$-dependent mapping of the deleted configuration. Its point measure is
the first component in \eqref{eq:extremes}. Its total weight is $n-J_1$,
and
\[
 \frac{J_1-D}{b}=-\frac{(n-J_1)-\mu_b}{b}.
\]
Apply Theorem~\ref{thm:cloud}. To pass from point-process convergence to
maxima, use \eqref{eq:largerpoints} to control points above an arbitrarily
large bound; the limiting process almost surely has no point at a
prescribed $x$. The number of limiting points above $x$ is Poisson with
mean $\int_x^\infty y^{-2}\dd y=1/x$. Zero such points gives
\eqref{eq:secondlargest}, and at most $k-1$ gives \eqref{eq:kthlargest}.
\end{proof}

\begin{corollary}[Separation of the condensate]\label{cor:separation}
In the dense regime $J_1/D\pto1$ and $J_2/J_1\pto0$. In the sparse regime
$J_1/n\pto1$ and again $J_2/J_1\pto0$.
\end{corollary}
\begin{proof}
In the dense case, $b/D\sim\ee\to0$, so tightness in
\eqref{eq:extremes} proves both claims. In the sparse case, bounded total
intensity and the uniform bound $\sum_{j>R}\nu_j\le C/R$ imply tightness
of the unconditioned largest action. Deletion in total variation transfers
this tightness to $J_2$. Also $K_M/n\pto0$ in the one-action proof gives
$J_1/n\pto1$. No assumption that $\lambda$ has a limit is needed for these
tightness statements.
\end{proof}

\section{Sharp action budgets and the logarithmic centering correction}
\label{sec:budgets}

\subsection{From a conditional maximum to a coefficient cutoff}
\begin{theorem}[Cutoff profile and its quantiles]\label{thm:budget}
If $\lambda\to\infty$, then \eqref{eq:maincutoff} and
\eqref{eq:mainbudget} hold for every real $x$ and every $r\in(0,1)$.
\end{theorem}
\begin{proof}
Combine the exact cutoff identity \eqref{eq:coefficientprob} with
Theorem~\ref{thm:extremes}. The floor changes the standardized threshold
by at most $1/b\to0$, and $F$ is continuous.
The threshold lies in $[0,n]$ eventually for every fixed $x$: $D/b\to\infty$,
and nonnegative rates and the exact tail imply
\[
 \frac{\mu_b}{b}\ge\sum_{J<j\le b}\frac1j\longrightarrow\infty,
\]
since $b^\ee j^{-\ee}\ge1$ for $j\le b$.
For a quantile, set $x_r=F^{-1}(r)$. For every $a>0$,
$F(x_r-a)<r<F(x_r+a)$. The corresponding cutoff ratios eventually bracket
$r$, so monotonicity brackets $(M_r-D)/b$ between $x_r-a+o(1)$ and
$x_r+a+o(1)$. Let $a\downarrow0$.
\end{proof}

The quantile law is asymptotic, not a finite-$n$ certified error bound.
In particular it does not license replacing $F$ by a floating-point
quadrature and calling the resulting integer cutoff a proof certificate.

\subsection{Why the leading coefficient scale is the wrong cutoff center}
\begin{theorem}[Logarithmically many fluctuation widths]\label{thm:centering}
In the dense regime,
\begin{align}
 \frac{D}{d}&\longrightarrow1,\label{eq:Doverd}\\
 \frac{D-d}{b}
 &=\frac{\log(1/\ee)}{1+\ee}
   +O\bigl(\ee\log^2(1/\ee)\bigr)+O(1/b).\label{eq:centering}
\end{align}
Thus the quantile formula has the explicit form
\begin{equation}\label{eq:explicitbudget}
 \boxed{\displaystyle
 M_r(n)=d_n+b_n\left(
       \frac{\log(1/\ee_n)}{1+\ee_n}+F^{-1}(r)\right)+o(b_n).}
\end{equation}
For every fixed real $C$,
\begin{equation}\label{eq:wrongcenter}
 \frac{u_{n,\lfloor d_n+Cb_n\rfloor}(\ee_n)}{u_n(\ee_n)}
 \longrightarrow0.
\end{equation}
\end{theorem}
\begin{proof}
From $D=b/\ee+O(1)$ and $d=b\ee^{-1/\alpha}$,
\[
 \frac{D}{d}=\ee^{-\ee/(1+\ee)}+O(1/d)\longrightarrow1.
\]
Furthermore
\[
 \frac{D-d}{b}
 =\frac1\ee\left(1-
       \exp\left(-\frac{\ee}{1+\ee}\log(1/\ee)\right)\right)+O(1/b).
\]
Taylor's formula for the ordinary exponential gives \eqref{eq:centering}.
The error there tends to zero, so Theorem~\ref{thm:budget} gives
\eqref{eq:explicitbudget}. For the last assertion,
$(d+Cb-D)/b\to-\infty$. Tightness of $(J_1-D)/b$ and the cutoff identity
then imply \eqref{eq:wrongcenter}.
\end{proof}

\begin{remark}[A strong failure of a leading-order prescription]
Both $d$ and $D$ are correct leading scales for the largest action.
Nevertheless a cutoff equal to $d$ plus any fixed number of cloud widths
retains an asymptotically vanishing fraction of the coefficient. This is
not a small error in a constant; it changes a desired fixed-retention
budget into a zero-retention prescription. The logarithm in
\eqref{eq:explicitbudget} is therefore an essential part of the endpoint
answer.
\end{remark}

\subsection{The macroscopic crossover}
\begin{corollary}[Condensate fraction at every logarithmic rate]
\label{cor:macro}
If $\ee_n\log n\to\tau\in[0,\infty]$, then
\begin{equation}\label{eq:macro}
 \frac{J_1}{n}\pto e^{-\tau},
\end{equation}
where $e^{-\infty}=0$. For any arbitrary-rate sequence and any cutoffs
with $M_n/d_n\to l\ne1$,
\begin{equation}\label{eq:macrocutoff}
 \frac{u_{n,M_n}}{u_n}\longrightarrow
 \begin{cases}0,&0\le l<1,\\1,&l>1.\end{cases}
\end{equation}
At $l=1$ no universal retained fraction is determined by this leading
ratio alone.
\end{corollary}
\begin{proof}
Directly from \eqref{eq:scales},
\begin{equation}\label{eq:logd}
 \log(d/n)=-\frac{\ee}{1+\ee}\log n
              +\frac1{1+\ee}\log(c_\ee/\ee).
\end{equation}
The last term tends to zero. The dense and sparse theorems give $J_1/d\pto1$;
for a sequence crossing between the two regimes, apply the subsequence
argument in Section~\ref{sec:matching}. Equation~\eqref{eq:logd} proves
\eqref{eq:macro}, and convergence in probability gives
\eqref{eq:macrocutoff}. Equations~\eqref{eq:maincutoff} and
\eqref{eq:wrongcenter} exhibit different limits with $M_n/d_n\to1$.
\end{proof}

\begin{example}[Three concrete joint limits]\label{ex:regimes}
For $\ee_n=n^{-1/2}$, $\lambda\sim\sqrt n\to\infty$ and the largest
action occupies asymptotically all the mass, while its cloud has scale
$b\sim\sqrt n$. A logarithmic budget shift of order $\sqrt n\log n$
is still present between $d$ and the correct center.

For $\ee_n=\tau/\log n$ with $0<\tau<\infty$,
\[
 d\sim e^{-\tau}n,\qquad
 b\sim\frac{\tau e^{-\tau}n}{\log n},\qquad
 D-d\sim\frac{\tau e^{-\tau}n\log\log n}{\log n}.
\]
Every row is exactly mean-critical, but the conditioned largest action
carries the positive fraction $e^{-\tau}$.

For $\ee_n=(\log n)^{-1/2}$, the largest action is sublinear in $n$,
but still asymptotically unique on its own scale $d$. It remains much
larger than the next action: $d/b=\ee^{-1/(1+\ee)}\to\infty$.
\end{example}

\section{Matching the coefficient laws and the critical transseries}
\label{sec:matching}

\subsection{The all-rate coefficient theorem}
\begin{proof}[Proof of Theorem~\ref{thm:main}]
First suppose $\lambda\to\infty$. Theorem~\ref{thm:dense} gives
\[
 \Prob(K=n)\sim\lambda D^{-1-\alpha}
   \sim\frac{\ee^{1+\alpha}}b.
\]
The proposed universal expression is
\[
 \frac\ee d=\frac{\ee^{1+1/\alpha}}b.
\]
Their ratio is $\ee^{\alpha-1/\alpha}\to1$, because
$(\alpha-1/\alpha)\log\ee\to0$. Combine this with
Proposition~\ref{prop:exact}.

If $\lambda=O(1)$, then $\ee\log n\to0$, $c_\ee/\ee\to1$, and
$d/n\to1$. Theorem~\ref{thm:sparse} gives
\[
 \frac{\Prob(K=n)}{\ee/d}
 \sim\frac{nc_\ee n^{-2-\ee}}{\ee/d}
 =\frac{c_\ee}{\ee}\frac dn n^{-\ee}\longrightarrow1.
\]
Deletion and separation have already been proved in both regimes.

For an arbitrary sequence, every subsequence has a further subsequence
on which $\lambda$ is bounded or tends to infinity: if it has a bounded
subsequence use it; otherwise select a subsequence diverging to infinity.
If any of the claimed scalar errors or total variation errors failed to
tend to zero, select a subsequence on which it stays bounded away from
zero and apply this dichotomy to obtain a contradiction. The same
argument applies to convergence-in-probability error probabilities.
Thus part (a) holds without restrictions on $\lambda$.
Part (b) is Theorems~\ref{thm:cloud}, \ref{thm:extremes}, and
\ref{thm:budget}, with \eqref{eq:Dproof}. Part (c) is
Theorem~\ref{thm:sparse}.
\end{proof}

\subsection{The familiar singular coefficient survives, but needs a new proof}
Set $p=1/(1+\ee)$ and recall $B_\ee=c_\ee\Gamma(-1-\ee)$.
The leading singular term in \eqref{eq:chart} is
$-B_\ee^{-p}h^p$. Its familiar coefficient expression is
\begin{equation}\label{eq:hahncoefficient}
 H_n(\ee)=\rho_\ee^{-n}
          \frac{B_\ee^{-p}}{-\Gamma(-p)}n^{-1-p}.
\end{equation}
It is positive for $0<p<1$. A fixed-$p$ singularity-transfer heuristic
suggests this expression, but is not a uniform argument as $p\to1$:
the reciprocal gamma factor vanishes at that endpoint.

\begin{corollary}[Uniform recovery of the singular coefficient]
\label{cor:hahncoefficient}
For every sequence $n\to\infty$, $\ee_n\to0^+$,
\begin{equation}\label{eq:hahnequiv}
 u_n(\ee_n)\sim H_n(\ee_n).
\end{equation}
More precisely the ratio of $H_n$ to the elementary right side of
\eqref{eq:maincoefficient} depends only on $\ee$:
\begin{equation}\label{eq:constantcorrection}
 C(\ee)=
 \bigl(\ee\Gamma(-1-\ee)\bigr)^{-1/(1+\ee)}
 \frac{-1}{\ee\Gamma(-1/(1+\ee))}
 =1-\ee+O(\ee^2).
\end{equation}
\end{corollary}
\begin{proof}
Cancellation of the powers of $n,c_\ee,\ee$ gives exactly
\eqref{eq:constantcorrection}. The first factor tends to one.
For the second, $1-p=\ee/(1+\ee)$ and
$-1/\Gamma(-p)\sim1-p\sim\ee$, so it also tends to one.
Combining with Theorem~\ref{thm:main} proves \eqref{eq:hahnequiv}.
For the displayed expansion, use
\[
 \ee\Gamma(-1-\ee)=1+(\gamma-1)\ee+O(\ee^2),
\]
and, for $v=1-p$, $-1/\Gamma(-1+v)=v+(\gamma-1)v^2+O(v^3)$.
The two first-order factors are
$1-(\gamma-1)\ee+O(\ee^2)$ and
$1+(\gamma-2)\ee+O(\ee^2)$, whose product is as stated.
\end{proof}

This corollary establishes a uniform coefficient statement without needing
a uniform global complex contour or an unproved cancellation of all
analytic terms at the moving endpoint. The local chart and the probability
proof describe the same leading coefficient by independent routes. The
corollary does \emph{not} give a relative error of order $o(\ee)$ for the
true coefficients; \eqref{eq:constantcorrection} is an exact comparison of
two candidate leading expressions, not a bound on their remaining errors.

\begin{remark}[Keep the exponential normalizer exact]
The factor $\rho_\ee^{-n}=\exp(nc_\ee A_\ee(1))$ is exact in all displayed
coefficient equivalents. Replacing its exponent by only the first term of
its small-$\ee$ expansion generally leaves an error of order $n\ee^2$
in the exponent. That error need not tend to zero in an arbitrary-rate
limit. The uniform theorem therefore does not authorize such a replacement
without an additional rate assumption. The same care is required when a
leading scale is used as a fluctuation center.
\end{remark}

\subsection{Matching sparse coefficients}
If $\lambda_n\to\lambda_0\in(0,\infty)$, then
\begin{equation}\label{eq:sparsecoefficient}
 u_n(\ee_n)\sim
 e^{\lambda_0(\zeta(2)+P(1))}\frac{\lambda_0}{n^3}.
\end{equation}
If $\lambda_n\to0$, the useful relative formula is instead
$u_n\sim\ee_n/n^2$. These follow from the sparse local theorem, since
$\rho_\ee^{-n}=\exp(\lambda A_\ee(1))$.
The endpoint coefficient for fixed $n$ from Proposition~\ref{prop:fixedq}
is consistent with this very sparse regime, but it does not by itself
justify letting $n$ grow.

\subsection{Why no contradiction with fixed-exponent behavior arises}
Fixing $\ee>0$ and letting $n\to\infty$ is not a sequence covered by the
endpoint theorem. In particular the assertion that all but one action
recover an unconditioned configuration in total variation is an endpoint
statement, not a claim for every fixed $1<\alpha<2$. The parameter
$\ee\to0$ makes the exceptional action larger than the ordinary cloud by
an unbounded factor. Similarly, finite-$\ee$ numerical ratios need not
approach one when compared with the elementary constant in
\eqref{eq:maincoefficient}; their limiting constant can differ by
$C(\ee)$. This distinction is visible in the diagnostics below.

\section{Computation, exact checks, and reproducibility}
\label{sec:computations}
The archive contains \texttt{verify.py}, its recorded
\texttt{verification\_results.json}, and the source of this article.
The checks are supplementary; no proof above relies on extrapolating a
numerical table.

\subsection{Exact algebra}
The script uses rational arithmetic in three different finite weight
models. It computes the formal solution of the feedback equation through
degree twelve and checks each coefficient against Lagrange inversion.
It independently enumerates multiplicity vectors through degree eight and
checks the exponential coefficient formula. These are 60 exact rational
identities. Two further symbolic identities check the first two equations
for the chart coefficients in Proposition~\ref{prop:chartcoeff}. All 62
exact assertions passed in the recorded run.

These tests validate the finite coefficient bookkeeping and the displayed
chart expansion. They are not a finite substitute for the local
probability estimates, which concern unbounded indices and moving
parameters.

\subsection{Positive coefficient recurrence}
For the numerical model, a useful recurrence avoids cancellation. If
$p_k=\Prob(K=k)$, then
\begin{equation}\label{eq:panjer}
 p_0=e^{-\sum_j\nu_j},\qquad
 p_k=\frac1k\sum_{j=1}^k j\nu_jp_{k-j}.
\end{equation}
This follows by differentiating the probability generating function. To
compute the joint probability of $K=n$ and no action above $M$, set the
rates above $M$ to zero in the recurrence but retain the \emph{original}
$p_0$. That convention preserves the probability of excluding the tail,
so division by the untruncated $p_n$ gives the exact mathematical cutoff
ratio. The floating implementation is not exact arithmetic.

The recorded run uses extended-precision \texttt{numpy.longdouble}
recurrences, with 70-digit \texttt{mpmath} normalization and chart
calculations. It checks that the initial and final probabilities are
positive and finite. Different platforms may provide a narrower
\texttt{longdouble}; underflow raises an error rather than silently
producing a result. The JSON records the mantissa precision.

\begin{table}[htbp]
\centering
\caption{Coefficient diagnostics for $P=0$. The elementary ratio is $p_n/(\epsilon/d)$; the gamma ratio uses $n\rho^nH_n$. Holding $\epsilon$ fixed is not the joint limit in the theorem.}
\label{tab:coefficients}
\begin{tabular}{rrrr}
\toprule $n$ & $\epsilon$ & Elementary ratio & Gamma ratio \\
\midrule
100 & 0.100 & 0.88565664 & 0.99693488 \\
1500 & 0.100 & 0.88808230 & 0.99966531 \\
4000 & 0.100 & 0.88824635 & 0.99984997 \\
100 & 0.020 & 0.97903537 & 0.99964242 \\
1500 & 0.020 & 0.97935966 & 0.99997353 \\
4000 & 0.020 & 0.97937548 & 0.99998969 \\
100 & 0.005 & 0.99488134 & 0.99992138 \\
1500 & 0.005 & 0.99495423 & 0.99999463 \\
4000 & 0.005 & 0.99495755 & 0.99999797 \\
\bottomrule
\end{tabular}
\end{table}

\begin{table}[htbp]
\centering
\caption{Genuine triangular diagnostics: $\epsilon_n=n^{-1/2}$ and a sparse sequence $\epsilon_n=1/n$. All entries use $P=0$ and the elementary normalization in the main theorem.}
\label{tab:triangular}
\begin{tabular}{rrrr}
\toprule $n$ & $\epsilon_n=n^{-1/2}$ & Dense ratio & Sparse ratio \\
\midrule
100 & 0.100000 & 0.88565664 & 0.98967840 \\
500 & 0.044721 & 0.95224627 & 0.99798733 \\
1500 & 0.025820 & 0.97314044 & 0.99933193 \\
4000 & 0.015811 & 0.98379152 & 0.99974980 \\
\bottomrule
\end{tabular}
\end{table}

\begin{table}[htbp]
\centering
\caption{Cutoff diagnostics, $P=0$, $\epsilon=0.02$, and $M=\lfloor D+bx\rfloor$. Fixed-$\epsilon$ bias remains; these values are not finite-$n$ error certificates.}
\label{tab:cutoff}
\begin{tabular}{rrrrr}
\toprule $n$ & $x$ & $M$ & Retained fraction & Limit $F(x)$ \\
\midrule
1500 & -4 & 1276 & 0.351031 & 0.296865 \\
1500 & -2 & 1332 & 0.526884 & 0.472186 \\
1500 & 0 & 1387 & 0.817415 & 0.789298 \\
1500 & 1 & 1415 & 0.951970 & 0.942065 \\
4000 & -4 & 3339 & 0.351476 & 0.296865 \\
4000 & -2 & 3484 & 0.525311 & 0.472186 \\
4000 & 0 & 3629 & 0.817311 & 0.789298 \\
4000 & 1 & 3702 & 0.951051 & 0.942065 \\
\bottomrule
\end{tabular}
\end{table}

\subsection{Stable quantiles and chart checks}
For the reflected law the inversion formula is
\begin{equation}\label{eq:cdfintegral}
 F(x)=\frac12+\frac1\pi\int_0^\infty
 e^{-\pi t/2}\frac{\sin(t(x+1-\gamma-\log t))}{t}\dd t.
\end{equation}
The apparent singularity at zero is integrable. The code uses numerical
quadrature on a finite interval with an exponentially small omitted tail;
its reported integration error is the quadrature routine's estimate, not
a rigorous enclosure. Representative values and finite-$n$ cutoff ratios
are in Table~\ref{tab:cutoff}. At these moderate parameters a visible
finite-$\ee$ bias remains. The limiting theorem assumes both
$\ee\to0$ and $b\to\infty$, rather than $n\to\infty$ with $\ee$ fixed.

For $\ee\in\{0.2,0.05,0.001\}$ and
$x\in\{0.01,0.001\}$, the program checks the exact phase against a
30-term analytic remainder and reconstructs $x$ by the contraction map
in the proof of Theorem~\ref{thm:chart}. The tested phase discrepancies
are below $10^{-60}$ and the reconstructed discrepancies below $10^{-55}$
in 70-digit arithmetic. Those thresholds describe the computation, not a
certified analytic remainder bound; the general bound is
\eqref{eq:chartbound}.

\subsection{What has not been computationally certified}
No interval-arithmetic cutoff algorithm is supplied. No numerical onset
for the asymptotic equivalences is claimed. No floating-point table is
called a proof of total variation convergence. No Lean file accompanies
this article. The package separates a reproducible mathematical proof,
exact finite algebra checks, and diagnostic floating calculations.

\section{A formalization plan with explicit proof boundaries}
\label{sec:formal}
The repository already has support and asymptotic-scale infrastructure
\cite{repo-project,repo-index}. One relevant module is
\path{TransseriesWellBased.lean}. That infrastructure should not be
mistaken for a formal proof of the new probability estimates. The natural
formalization is modular.

\paragraph{Finite algebra and exact coefficients.}
First prove the exponential-series recurrence over a rational algebra and
its multiplicity-vector formula. Connect the formal implicit solution to
Lagrange inversion. The positivity and finite-cutoff monotonicity
statements can then be developed without analytic limits. A certificate
checker for a finite list of rational weights should sit above these
proved identities, not replace them.

\paragraph{Finite and countable Poisson configurations.}
Construct finite-support count vectors under a summable rate sequence,
prove the generating-function identity, and connect the cutoff event to
the coefficient ratio. The convolution estimate for multiple large
actions is an $\ell^1$--$\ell^\infty$ inequality and is a useful early
formal target. The deletion map is elementary on the count-vector space.

\paragraph{Uniform local control.}
Formalize the integral-test bounds with explicit constants on
$1\le\alpha\le5/4$, then the trigonometric lower bound and Fourier
concentration estimate. The exponential-tilt identity should expose its
hypotheses, particularly positivity of the tilt and the finite cutoff.
Only then assemble the zero/one/multiple-action proof. The exact
$L^1$ density argument is the central bridge from scalar estimates to
total variation.

\paragraph{Moving-parameter convergence.}
A reusable Riemann-sum theorem with a parameter in a compact interval,
a small-jump variance estimate, and a truncation lemma for joint weak
convergence would support Theorem~\ref{thm:cloud}. Quantile convergence
requires a separate continuity and strict-monotonicity interface; it is
not an automatic consequence of pointwise coefficient asymptotics.

\paragraph{Analytic chart.}
The required ingredients are the polylogarithm expansion, cancellation of
meromorphic poles in the parameter, and a parameter-uniform holomorphic
implicit-function or contraction theorem. Ordinary multivariable power
series can be used before constructing the specialized Hahn support.
This keeps analytic existence distinct from formal support manipulation.

This is a plan, not a claim that these interfaces are already implemented
in the pinned repository. In particular the present article adds no
formal declarations to that repository and makes no changes to its files.

\section{Further research questions}
\label{sec:future}
The results above answer the exact-tail lower-endpoint question at leading
coefficient order and at the sharp fluctuation and cutoff scale. They also
make the remaining questions more specific.

\paragraph{1. Effective error bounds and certified action selection.}
The proof gives explicit inequalities before several asymptotic
simplifications. Optimize the splitting threshold $M=\eta D$, retain
constants in the Fourier concentration bound, and obtain a computable
upper bound on the deletion distance. Combine that bound with rigorous
enclosures of \eqref{eq:cdfintegral} to return a certified finite-$n$
cutoff for a prescribed retained fraction. The difficult step is not
numerical integration alone; it is a quantitative comparison between the
finite conditional law and its limit.

\paragraph{2. A uniform second-order coefficient theorem.}
Equation~\eqref{eq:constantcorrection} compares two leading normalizations
but does not bound the remainder by $o(\ee)$. Determine a joint expansion
in $\ee$, $1/b$, and the logarithmic combinations generated by
$\ee\log\ee$. A useful target is an error formula valid across
$\lambda=O(1)$, $\lambda\to\infty$ slowly, and very large $b$, rather than
three unrelated expansions with no common overlap.

\paragraph{3. Slowly varying and nonexact tails.}
Replace $j^{-2-\ee}$ by $j^{-2-\ee}L_\ee(j)$. Identify uniform regularity
conditions on $L_\ee$ that preserve one-action condensation, and construct
the correct cloud and mean-defect scales. Ordinary regular variation for
each fixed parameter need not be uniform as the parameter moves. Determine
when a de Bruijn conjugate or a more general implicit scale replaces the
explicit powers $b$ and $d$.

\paragraph{4. Simultaneous coupling displacement.}
Allow the mean of $K$ to be $n(1-\delta_n)$ rather than $n$ while
$\ee_n\to0$. Find the coupling window in which the imposed deficit,
truncated-mean defect, and Cauchy cloud all interact. The natural
centering candidate is $n-\sum_{j\le b}j\nu_j$, but positivity of that
quantity and local one-action dominance must be re-established. The
present proof does not address the side on which conditioning lies below
the typical cloud.

\paragraph{5. A growing prefix and the loss of universality.}
Let the altered prefix length $J_n$ grow. Determine sharp conditions on its
centered variance, maximum action, and tail intensity under which it is
still invisible after scaling by $b_n$. A prefix with negligible total
mass can nevertheless shift a fluctuation center if its first moment is
large. The fixed-prefix proof isolates exactly where stronger conditions
would be required.

\paragraph{6. Several action species and competing condensates.}
For a sum of several near-index-one tails, mark the species of each action.
Determine when exactly one condensate still occurs and how its species is
selected. If different exponents approach one at different rates, the
largest-action weight at the mean defect may select one species while the
ordinary cloud contains several. Prove the corresponding marked
point-process and cutoff statements before assuming a single effective
exponent.

\paragraph{7. Geometry beyond the coefficient configuration.}
The conditioned Poisson configuration encodes coefficient contributions,
not by itself a complete tree metric. Use an explicit combinatorial
bijection to connect this triangular model to tree degree sequences, and
then investigate height, diameter, and scaling limits. The fixed-law
Cauchy-tree literature \cite{kr} supplies relevant comparisons, but its
geometric conclusions cannot be transferred without that additional
structure.

\paragraph{8. Complex continuation and additional transseries structure.}
Theorem~\ref{thm:chart} is a convergent local critical chart. It does not
establish global analytic continuation, resurgence, Borel summability, or
Stokes constants. Determine what happens when $h$ and $\ee$ vary in larger
complex domains, when the chosen branch winds, or when the parameter
approaches other singularities of the normalization. State separately
which properties follow from a finite-coordinate analytic chart and
which require new global information.

\paragraph{9. Machine-checked deletion and triangular limits.}
Formalize the deletion theorem as a reusable interface: a negligible
zero-large-action term, a negligible multi-large-action term, and an
$L^1$-asymptotically constant one-action weight imply total variation
recovery of the cloud. This abstract statement would support many
condensation models beyond the explicit zeta family, while leaving the
model-specific tail and concentration estimates as separate obligations.

\section{Conclusion}
The lower endpoint does not merely replace one stable exponent by another.
It turns the vanishing critical coupling into a mean-defect problem, with
one exceptional action and a smaller compensated Cauchy cloud. For the
exact power tail with a fixed finite prefix, the coefficient equivalent
holds at every rate of approach, including regimes in which the total
Poisson intensity stays bounded or tends to zero. In the dense regime the
cutoff quantiles require a logarithmic correction to the leading scale;
omitting it gives asymptotically zero retention even though the cutoff is
correct to relative leading order.

The analytic and probabilistic descriptions fit together without a hidden
uniformity assumption. A convergent critical chart persists through a
removable parameter singularity, while a separate local probability proof
establishes the uniform coefficient law. The matching gamma expression is
therefore a conclusion, not a substitution-based justification. The
remaining tasks concern quantitative certification, less rigid tails,
additional moving parameters, and formal verification.

\appendix
\section{Proof-dependency and scope ledger}
\label{app:ledger}
\begin{longtable}{@{}p{0.28\textwidth}p{0.64\textwidth}@{}}
\toprule
Statement & Proof inputs and boundary \\
\midrule\endhead
Exact coefficients and cutoffs & Formal Lagrange inversion and summable
Poisson rates; reproved in Proposition~\ref{prop:exact}. No asymptotic
assumption.\\[5pt]
Uniform inverse chart & Classical polylogarithm expansion, removable pole
cancellation, and a local holomorphic implicit-function argument. Local
branches only.\\[5pt]
Dense local equivalent & Integral-test estimates, a tilted Fourier
concentration bound, and a decomposition into zero, one, or multiple
large actions. No fixed-exponent local limit is imported.\\[5pt]
Sparse local equivalent & A positive exponential-moment bound retaining a
factor of $\lambda$, plus the same one-action decomposition. Valid when
$\lambda$ tends to zero.\\[5pt]
Deletion in total variation & An explicit bounded Radon--Nikodym weight
converging to one in $L^1$; no moment convergence asserted.\\[5pt]
Cauchy cloud & Poisson Laplace functionals, compact Riemann sums, uniform
upper-tail bounds, and small-jump variance control. Requires $b\to\infty$.\\[5pt]
Sharp action budgets & Exact cutoff identity, deletion, weak convergence,
and strict continuity of the limiting quantile. Not a finite-$n$
certificate.\\[5pt]
All-rate coefficient law & Dense and sparse local equivalents joined by a
subsequence dichotomy; fixed prefix throughout.\\[5pt]
Publication novelty & Not established. The arbitrary-rate theorem is
presented as a proved, scoped contribution; classical condensation
mechanisms are credited.\\[5pt]
Lean status & No new Lean proofs or repository modifications are included.\\
\bottomrule
\end{longtable}

\section{Reproduction instructions and source provenance}
The article is self-contained at the level of the displayed conventional
proofs. Its external mathematical input is identified in the bibliography.
Repository claims are tied to \commit, not to a later moving
\texttt{main} branch. The source notes in the archive list the inspected
paths and the specific question answered.

To rebuild the PDF from this directory, run
\begin{verbatim}
sh build.sh
\end{verbatim}
To reproduce all recorded checks, install the packages in
\texttt{requirements.txt} and run
\begin{verbatim}
python verify.py
\end{verbatim}
The optional \texttt{--quick} run writes a separate file,
\texttt{verification\_quick.json}, so it does not replace the full recorded
results. Floating results can differ in their last digits across library
versions and hardware. The program checks failure conditions but does not
perform directed rounding.

\begin{thebibliography}{9}
\bibitem{repo-critical}
V.~Reshetnikov, \emph{ProveIt repository: Beyond Finite-Action Folds:
Critical Hahn Transseries, Stable Sector Asymptotics, and Sharp Action
Budgets}, research report, 29 September 2026.
\href{https://github.com/VladimirReshetnikov/ProveIt/blob/3d5973524506411392a911470b5ddc35521568ea/Analysis/Transseries/docs/series-and-transseries/Critical_Hahn_Transseries_Beyond_Finite_Action_Folds/article.tex}{Pinned TeX source}.
Question~4, source lines 1502--1507, and the adjacent editorial note;
model and exact coefficient identities in the opening sections.
Repository ownership is recorded here, not an assertion about individual
human authorship of every research draft.

\bibitem{repo-project}
V.~Reshetnikov, \emph{ProveIt: Analysis/Transseries}, project README at
revision \texttt{3d5973524506}.
\href{https://github.com/VladimirReshetnikov/ProveIt/blob/3d5973524506411392a911470b5ddc35521568ea/Analysis/Transseries/README.md}{Pinned project description}.

\bibitem{repo-index}
V.~Reshetnikov, \emph{ProveIt: Series and transseries}, documentation
index at revision \texttt{3d5973524506}.
\href{https://github.com/VladimirReshetnikov/ProveIt/blob/3d5973524506411392a911470b5ddc35521568ea/Analysis/Transseries/docs/series-and-transseries/README.md}{Pinned documentation index}.
In particular the editorial accounts of the upper-endpoint research
packages and their remaining lower-endpoint question.

\bibitem{janson}
S.~Janson, \emph{Simply generated trees, conditioned Galton--Watson trees,
random allocations and condensation}, Probability Surveys \textbf{9}
(2012), 103--252. \href{https://doi.org/10.1214/11-PS188}{doi:10.1214/11-PS188}.
\href{https://arxiv.org/abs/1112.0510}{arXiv:1112.0510}.

\bibitem{al}
I.~Armend\'ariz and M.~Loulakis, \emph{Conditional distribution of heavy
tailed random variables on large deviations of their sum}, Stochastic
Processes and their Applications \textbf{121}(5) (2011), 1138--1147.
\href{https://doi.org/10.1016/j.spa.2011.01.011}{doi:10.1016/j.spa.2011.01.011}.
\href{https://arxiv.org/abs/0912.1516}{arXiv:0912.1516}.

\bibitem{kr}
I.~Kortchemski and L.~Richier, \emph{Condensation in critical Cauchy
Bienaym\'e--Galton--Watson trees},
\href{https://arxiv.org/abs/1804.10183v3}{arXiv:1804.10183v3},
21 November 2018. The cited version is the authors' final preprint;
see its condensation and conditioned-walk limit theorems.

\bibitem{dlmf}
NIST Digital Library of Mathematical Functions,
\emph{Section 25.12: Polylogarithms}, especially equation 25.12.12.
\url{https://dlmf.nist.gov/25.12.E12}. Accessed 29 September 2026.
\end{thebibliography}
\end{document}
